\begin{abstract}
Frequent pattern mining identifies itemsets whose co-occurrence frequency exceeds a support threshold, and sliding-window extensions can detect temporal intervals where this frequency is locally elevated.
However, attributing such support changes to specific external events---e.g., a marketing campaign or policy intervention---remains an open problem: observed correlations may be spurious.
We propose \textsc{SC-DenseAttrib}, the first method to apply the Synthetic Control Method (SCM) to pattern support time series for causal attribution.
Our key idea is to construct a \emph{donor pool} of item-disjoint control patterns whose support trajectories are unaffected by the event of interest, then estimate a counterfactual support trajectory as a weighted combination of these controls.
The causal effect is measured as the divergence between observed and counterfactual support after the intervention.
We provide formal definitions of Donor Pool, Counterfactual Support Trajectory, Causal Effect on Support, and Item-Disjoint Pattern Control, and prove unbiasedness under parallel trends.
Experiments on synthetic data show that SC-DenseAttrib recovers injected causal effects with mean absolute error below 0.04, maintains a false positive rate of 0\% under the null, and provides interpretable counterfactual trajectories that simple permutation tests cannot offer.
\end{abstract}
